\chapter{63}



Wider Erwarten scharf erklang der erste Klingelton. Louis’ Körper zog
sich zusammen. Noch konnte er den Anruf unerkannt beenden. Das Klingeln
wiederholte sich, drei, vier, fünfmal, dann brach der Kontakt ab. Da war
es wieder, dieses lähmende Gefühl. Doch was sollte denn schon geschehen?
Was immer er auch von Valentina erfahren würde, schlimmer konnte es
nicht werden. Im Gegenteil. Es würde hilfreich sein, zu wissen, was sich
gerade in Mailand abspielte.

\sectionbreak

Beim erneuten Eingeben der Nummer vertippte sich Louis. Zwei Mal. Das dritte Mal klappte es, die Verbindung stand. 
Doch wieder war nach fünf Mal klingeln Schluss. Er liess das Gerät auf seinen Schoss fallen.
Dann veränderte sich alles. Jetzt schrillte \emph{sein} Handy. Louis
starrte auf das Display und sah die Nummer. Wie war das möglich, hatte
er vergessen -, seine Gedanken überschlugen sich und wie in Trance drückte
er auf den grünen Hörerknopf. Erst blieb es auf der andern Seite still.

Dann sprach Valentina. Ihre Worte klangen bestimmt. Im Hintergrund war
ein Durcheinander an Stimmen zu hören.

«Wer ist da?»

Selbst jetzt konnte er den Anruf noch beenden. Doch was nützte es. Sie
hatte die Nummer. Er hatte tatsächlich vergessen, die Einstellung auf
unterdrückte Nummer zu wechseln. Wie dumm war er eigentlich. Sein Gehirn
ratterte auf Hochtouren. Er musste das Beste daraus machen. An ihre
Freundschaft appellieren und auf ihre Verschwiegenheit hoffen.

«Ich bin`s, Louis.»

«Louis?»

Er hörte Valentina atmen.

«Ja, hast du einen Moment? Bist du alleine?»

Die zweite Frage war ihm rausgerutscht. Valentina sprach gedämpft.

«Ich bin gerade auf der Piazza, wir sind beim Essen, warte, ich setze
mich mal ab, Moment.»

Die Stimmen wurden leiser und verschwanden ganz.

«Louis, ich glaub es ja nicht. Wo bist du?»
